\documentclass[11pt]{article}
\usepackage[a4paper,margin=27mm]{geometry}
\usepackage{amsmath,amssymb,mathtools}
\usepackage{fontspec,unicode-math}
\setmainfont{Latin Modern Roman}
\setmathfont{Latin Modern Math}
\usepackage{xurl}
\usepackage[hidelinks]{hyperref}
\newcommand{\sref}[2]{\hyperlink{source:#1:#2}{[S#2]}}
\newcommand{\cataloguescope}[1]{\paragraph{Mathematical question.}#1}
\hypersetup{pdftitle={UnsolvedMath Problem 9400095}}
\begin{document}
% UnsolvedMath ID 9400095; source catalogue code AMR-093-0095
\setcounter{section}{9400094}
\section{Does every rational number with an odd denominator have an odd greedy expansion}\label{9400095}
{\small\sffamily UnsolvedMath ID 9400095 · Number Theory}
\subsection{Definitions and mathematical statement}

\paragraph{Shared notation.}$\mathbb N=\{1,2,\ldots\}$, and $\mathbb Z,\mathbb Q,\mathbb R,\mathbb C$ denote the integers, rationals, reals and complex numbers. Any other conventions are specified locally.


Write a positive rational number in lowest terms as $r_0=a/b$ with $a,b\in\mathbb N$ and $b$ odd. As long as $r_j>0$, choose the least positive odd integer $d_j$ with $1/d_j\le r_j$ and set $r_{j+1}=r_j-1/d_j$. The choice permits $d_j=1$ and repeated denominators. It is the largest available unit fraction not exceeding the remainder. The expansion terminates if some remainder becomes zero.
\paragraph{Statement recorded in the supplied source.}
Does this odd greedy algorithm terminate for every such $a/b$, yielding $a/b=\sum_{j=0}^{m-1}1/d_j$ with finitely many terms? \sref{9400095}{2}
\subsection{Short English statement}
Must the greedy algorithm restricted to odd denominators eventually express every positive rational with odd denominator as a finite sum of unit fractions?
\subsection{Sources}
\begin{description}
\item[\hypertarget{source:9400095:1}{S1}] ULAM.AI and the UnsolvedMath Contributors, UnsolvedMath: A Curated Collection of Open Mathematics Problems, record 9400095 (AMR-093-0095). CC BY 4.0. \url{https://huggingface.co/datasets/ulamai/UnsolvedMath}.
\item[\hypertarget{source:9400095:2}{S2}] Joel Louwsma and Joseph Martino, \emph{Rational numbers with odd greedy expansion of fixed length}, arXiv:2309.07280v1 (2023) Introduction, pp. 1--2. \url{https://arxiv.org/pdf/2309.07280}.
\item[\hypertarget{source:9400095:3}{S3}] Wikipedia contributors, \emph{List of unsolved problems in mathematics}, revision 1366636767 (2026) Number theory, Does every rational number with an odd denominator have an odd greedy expansion. \url{https://en.wikipedia.org/w/index.php?title=List_of_unsolved_problems_in_mathematics&oldid=1366636767#Number_theory}.
\end{description}
\label{end:9400095}
\end{document}
